\bookchapter{第二十九章：数字健康——和手机建立健康关系}

\chapterintro{你有没有这样的体验：本来只想刷5分钟短视频，结果一抬头已经过去了两小时？或者晚上躺在床上刷手机，越刷越精神，最后凌晨两点才睡着？手机改变了我们的生活，但也在悄悄吞噬我们的时间和健康。这一章不是让你“戒掉手机”，而是帮你重新夺回对手机的主动权。}

\booksection{一、先看看你的“屏幕时间”——数字让人清醒}

在改变之前，先搞清楚现状。手机里有一个很多人从未打开过的功能，叫“屏幕使用时间”（iPhone）或“数字健康”（安卓），它能告诉你每天到底在手机上花了多少时间，以及这些时间花在了哪些App上。

\begin{itemize}
  \item \textbf{iPhone}：设置 > 屏幕使用时间 > 查看所有活动。你会看到每天的解锁次数、拿起手机的次数、每个App的使用时长。
  \item \textbf{安卓}：设置 > 数字健康与家长控制。华为叫"健康使用手机"，小米叫"屏幕时间管理"。功能类似，都能看到每日使用数据。
\end{itemize}

\textbf{第一次看这个数据的人，80\%会被吓到}。我有个朋友看完之后发现自己每天在抖音上花了4个多小时，一周就是28小时——相当于一份兼职的时间。这种“数字冲击”是改变的第一步。

\booksection{二、给App设置“使用限额”——给自己装个限速器}

知道问题在哪之后，就可以有针对性地设置限制。不需要一刀切把App删了，只需要给它设个“每日额度”：

\begin{enumerate}
  \item \textbf{iPhone}：设置 > 屏幕使用时间 > App限额 > 添加限额。选中你要限制的App（比如抖音），设置每天30分钟。到时间后App图标会变灰，打开需要额外输入密码。
  \item \textbf{安卓}：设置 > 数字健康 > 应用计时器。选中App后设置每日使用上限，到时间后App图标会变灰。
  \item \textbf{小技巧}：把限额密码设成一个你不常用的、不太好记的密码。太容易输的话，你会忍不住每次都"续命"，限额就形同虚设了。
\end{enumerate}

\booksection{三、“勿扰模式”的正确用法——不是关掉通知，是关掉焦虑}

很多人的手机一天响几十次通知，每响一次注意力就被打断一次。你以为你只是“看一眼通知”，实际上每次打断后重新集中注意力需要15-20分钟。勿扰模式不是让你“失联”，而是帮你重新掌控注意力。

\begin{itemize}
  \item \textbf{定时开启}：设置勿扰模式在晚上10点到早上7点自动开启。这样睡觉时手机不会因为一条垃圾短信把你吵醒。
  \item \textbf{允许重要联系穿透}：在勿扰模式设置里，把"允许来电"设为"收藏联系人"或"重复来电"（3分钟内第二次来电会正常响铃）。这样你不会错过家人的紧急电话，但也不会被推销电话打扰。
  \item \textbf{工作专注模式}：iPhone的"专注模式"可以设自定义场景——比如"工作"模式下只允许工作相关App通知，"阅读"模式下所有通知全静音。
\end{itemize}

\booksection{四、护眼不靠App靠习惯——20-20-20法则}

长时间看手机屏幕，眼睛干涩、酸胀是必然的。护眼模式（蓝光过滤）有一定帮助，但真正有效的是改变用眼习惯：

\begin{tcbinfo}{20-20-20法则}
每看屏幕\textbf{20分钟}，把视线移到\textbf{20英尺（约6米）}以外的地方，保持\textbf{20秒}。这个简单的动作能有效缓解睫状肌疲劳。

\end{tcbinfo}

\begin{itemize}
  \item \textbf{开启护眼模式}：设置 > 显示 > 护眼模式（或夜览模式），开启后屏幕色温变暖，减少蓝光。建议设为日落到日出自动开启。
  \item \textbf{调整屏幕亮度}：环境暗时屏幕也要暗。很多人在黑暗中看手机亮度调太高，这是伤眼最严重的行为。开启"自动亮度"并确保它正常工作。
  \item \textbf{保持距离}：手机距离眼睛至少30厘米。很多人喜欢把手机拿得特别近看，这样睫状肌一直在紧张收缩状态，近视加深很快。
\end{itemize}

\booksection{五、“低头族”的颈椎自救指南}

低头看手机时，你的脖子承受的压力是正常站立时的5-6倍。长期低头会导致颈椎曲度变直、骨质增生、头晕手麻。以下是几个简单的改善方法：

\begin{itemize}
  \item \textbf{把手机举高}：看手机时尽量把手机举到与眼睛平齐的高度，而不是低头看放在腿上的手机。虽然手臂会累，但脖子会感谢你。
  \item \textbf{每30分钟活动一下脖子}：慢慢向左转头到极限，停3秒，回正；再向右转，停3秒。重复5次。然后做"米字操"——用下巴在空中写一个"米"字。
  \item \textbf{靠墙站立3分钟}：后脑勺、肩胛骨、臀部、脚后跟四点贴墙，每天站3分钟，能有效纠正含胸低头的体态。
\end{itemize}

\booksection{六、睡眠与手机——卧室里最该改的习惯}

手机屏幕的蓝光会抑制褪黑素分泌，让你“越看越精神”。很多人觉得“我刷着刷着就困了”，实际上是被蓝光欺骗了大脑，等你放下手机反而不容易入睡。以下是改善建议：

\begin{itemize}
  \item \textbf{睡前1小时不看手机}：这是最有效也最难做到的一条。如果做不到，至少开"夜览/护眼模式"把蓝光降到最低。
  \item \textbf{手机放在伸手够不到的地方}：把充电器放在卧室另一头，而不是床头柜。这样你既不会忍不住拿起来刷，闹钟响了你也得起床才能关掉——一举两得。
  \item \textbf{开启"睡眠模式"}：iPhone的"睡眠"专注模式可以自动在睡眠时段锁定屏幕、简化锁屏；安卓各品牌也有类似功能。
  \item \textbf{用闹钟代替手机}：买一个十几块钱的实体闹钟，不用手机当闹钟，你会发现少了"看一眼时间然后忍不住刷一下"的陷阱。
\end{itemize}

\booksection{七、信息焦虑症——学会“选择性失明”}

我们每天接收的信息量，是古人一辈子接收的量。但人的大脑不是为处理这么多信息设计的，信息过载会导致焦虑、注意力涣散、决策疲劳。解决方法不是“获取更多信息”，而是“主动屏蔽”：

\begin{itemize}
  \item \textbf{取消关注}：打开你关注的所有公众号、UP主、博主列表，问自己一个问题："过去一个月我从中获得过实际价值吗？"如果没有，取消关注。信息渠道贵精不贵多。
  \item \textbf{关闭非必要App推送}：除了微信、电话、短信、日历之外，其他App的通知全关掉。新闻不需要推送——你主动去看的时候自然会看到。
  \item \textbf{设立"信息斋戒日"}：每周选一天（比如周日），全天不看手机新闻、不刷短视频。你会发现天不会塌下来，而且周一精神状态明显更好。
\end{itemize}

\begin{tcbpractice}{动手：一周数字健康挑战}
\begin{enumerate}
  \item 今天：打开屏幕使用时间/数字健康，查看你上周的使用数据。
  \item 明天：给耗时最长的娱乐App设置每日限额（建议30-60分钟）。
  \item 后天：设置勿扰模式定时开启（22:00-07:00），允许收藏联系人来电。
  \item 第四天：睡前1小时不看手机，手机充电器放到卧室另一头。
  \item 第五天：取消关注5个你不再看的公众号或博主。
  \item 第六天：尝试一整天的"信息斋戒"——不看新闻不刷短视频。
  \item 第七天：查看本周的屏幕使用时间报告，和上周对比。
\end{enumerate}

\end{tcbpractice}

\begin{tcbsummary}{本章小结}
\begin{itemize}
  \item 先看屏幕使用时间数据，认清现状是改变的第一步
  \item 给娱乐App设限额，给勿扰模式设定时，重新掌控注意力
  \item 20-20-20法则是最简单有效的护眼习惯
  \item 低头看手机时脖子承重5-6倍，举高手机+定期活动脖子
  \item 睡前1小时不看手机，充电器放到伸手够不到的地方
  \item 信息焦虑的解药是"主动屏蔽"，不是"获取更多"
\end{itemize}

\end{tcbsummary}
